\section{Technical Approach}

\subsection{Pipeline}

Let $\mathcal{C}$ be a bank of English activity sentences, each associated
with a text embedding $e_c$. For every captured frame $x_i$, a lightweight
encoder running locally on the iPhone produces an image representation and
selects the draft caption
\begin{equation}
c_{d,i}
=
\arg\max_{c\in\mathcal{C}}
\cos\big(f_I(x_i),e_c\big).
\end{equation}

Because this draft stage is executed entirely on the iPhone, it does not
require wireless communication. The selected sentence $c_{d,i}$ can therefore
be displayed immediately after the local inference completes.

Every $\gamma$ drafts, or whenever another verification rule selects a frame,
the corresponding image is submitted for remote verification. The iPhone
first generates the JPEG representation to be uploaded and sends a small
verification request containing the image size and the desired end-to-end
confirmation deadline.

After receiving permission from the communication scheduler, the iPhone
uploads the image to the laptop edge server. SmolVLM-256M then generates a
short action phrase $s_i$. The generated phrase is mapped back to the same
sentence bank as
\begin{equation}
c_i^\star
=
\arg\max_{c\in\mathcal{C}}
\cos\big(f_T(s_i),e_c\big).
\end{equation}

The local draft is accepted when
\begin{equation}
c_i^\star=c_{d,i}
\end{equation}
and
\begin{equation}
\cos(e_{c_{d,i}},f_T(s_i))\geq\tau.
\end{equation}

Otherwise, the displayed caption is replaced by $c_i^\star$.

Therefore, only the verification path requires wireless communication:
\[
\mathrm{iPhone}
\rightarrow
\mathrm{Wi\mbox{-}Fi}
\rightarrow
\mathrm{Laptop~VLM}
\rightarrow
\mathrm{iPhone}.
\]

\begin{algorithm}[t]
\caption{Local draft and edge verification}
\label{alg:verify}
\begin{algorithmic}[1]
\STATE Run the lightweight local model on frame $x_i$.
\STATE $c_{d,i}\leftarrow
\arg\max_{c\in\mathcal{C}}
\cos(f_I(x_i),e_c)$.
\STATE Display $c_{d,i}$ immediately.
\IF{frame $i$ is selected for verification}
    \STATE Predict its edge verification latency $\hat T_i^{\mathrm{edge}}$.
    \STATE Compute its communication budget $B_i^{\mathrm{up}}$.
    \STATE Submit the verification request to the communication scheduler.
    \STATE Upload the frame when scheduled.
    \STATE Generate phrase $s_i$ with SmolVLM.
    \STATE $c_i^\star\leftarrow
    \arg\max_{c\in\mathcal{C}}
    \cos(f_T(s_i),e_c)$.
    \IF{$c_i^\star=c_{d,i}$ \AND
    $\cos(e_{c_{d,i}},f_T(s_i))\geq\tau$}
        \STATE Accept $c_{d,i}$.
    \ELSE
        \STATE Replace the caption with $c_i^\star$.
    \ENDIF
\ENDIF
\end{algorithmic}
\end{algorithm}

\subsection{Latency Model}

Two latency paths are considered separately.

The draft latency is purely local:
\begin{equation}
T_{i}^{\mathrm{draft}}
=
T_{i}^{\mathrm{local}},
\end{equation}
where $T_{i}^{\mathrm{local}}$ is the execution time of the lightweight model
on the iPhone. No wireless latency is involved in the draft path.

For a frame selected for VLM verification, the confirmation delay is
\begin{equation}
T_i^{\mathrm{conf}}
=
T_i^{\mathrm{wait}}
+
T_i^{\mathrm{up}}
+
T_i^{\mathrm{edge}}
+
T_i^{\mathrm{down}},
\end{equation}
where $T_i^{\mathrm{wait}}$ is the waiting time before the communication
scheduler permits the upload, $T_i^{\mathrm{up}}$ is the image upload time,
$T_i^{\mathrm{edge}}$ is the edge-side VLM verification latency, and
$T_i^{\mathrm{down}}$ is the return delay of the verification result.

The downlink payload is small because the server only returns an acceptance
decision or a caption identifier. Therefore, the dominant components are
typically the uplink transmission and VLM computation.

Each verification request $i$ is associated with an end-to-end confirmation
deadline
\[
D_i^{\mathrm{conf}}.
\]
The request is considered successful if
\begin{equation}
T_i^{\mathrm{conf}}
\leq
D_i^{\mathrm{conf}}.
\end{equation}

\subsection{Online Edge-Latency Prediction}

The execution time of the VLM is not assumed to be constant. It may change
with the uploaded image and the current computational state of the laptop.
We therefore use a lightweight predictor to estimate the edge-side
verification latency before the full image is uploaded.

For verification request $i$, the predictor uses
\begin{equation}
\mathbf{z}_i
=
\left[
L_i,\;
U_i^{\mathrm{CPU}},\;
T_{i-1}^{\mathrm{VLM}}
\right],
\end{equation}
where $L_i$ is the encoded JPEG size,
$U_i^{\mathrm{CPU}}$ is the current CPU utilization of the laptop, and
$T_{i-1}^{\mathrm{VLM}}$ is the measured VLM latency of the previous
verification request.

A small neural network $g_\theta(\cdot)$ predicts
\begin{equation}
\hat T_i^{\mathrm{edge}}
=
g_\theta(\mathbf{z}_i).
\end{equation}

The predictor will be trained from latency traces collected on the actual
laptop. The model is intentionally kept small because latency prediction is
a supporting component rather than the main AI workload of the system.

Prediction accuracy will be measured using mean absolute error between the
predicted and measured VLM latency.

\subsection{Compute-Aware Link Control}

The system exposes three communication-related control mechanisms.

First, the verification period $\gamma$ determines how frequently frames are
sent to the edge VLM. Second, JPEG quality and resolution determine the number
of bytes that must be uploaded. These two controls determine the verification
load and communication volume.

Third, and independently of these application-level controls, we introduce a
compute-aware upload scheduler for frames that have already been selected for
verification.

For request $i$, after predicting the edge latency
$\hat T_i^{\mathrm{edge}}$, the remaining uplink budget is
\begin{equation}
B_i^{\mathrm{up}}
=
\left[
D_i^{\mathrm{conf}}
-
\hat T_i^{\mathrm{edge}}
-
\hat T_i^{\mathrm{down}}
-
M
\right]^+,
\end{equation}
where $\hat T_i^{\mathrm{down}}$ is the expected result-return latency and
$M$ is a safety margin used to tolerate prediction error.

If request $i$ enters the communication scheduler at time $a_i$, its derived
uplink completion deadline is
\begin{equation}
d_i^{\mathrm{up}}
=
a_i+B_i^{\mathrm{up}}.
\end{equation}

With multiple phones waiting to upload verification frames, the scheduler
prioritizes the request with the earliest derived uplink deadline:
\begin{equation}
i^\star
=
\arg\min_{i\in\mathcal{Q}}
d_i^{\mathrm{up}},
\end{equation}
where $\mathcal{Q}$ is the set of pending verification uploads.

This scheduling mechanism is implemented above the Wi-Fi MAC layer. We do not
assume control of physical frequency resources. Instead, each client first
sends a small verification request and waits for an application-layer upload
grant from the laptop. The selected client then uploads its JPEG frame.

The key idea is that wireless urgency depends on downstream computation. If
the edge server is predicted to require a long time to verify a frame, less
time remains for communication and that request receives an earlier uplink
deadline. When the predicted edge latency is short, the communication
scheduler can safely tolerate a longer waiting time.